\documentclass[11pt]{article}
\usepackage[T1]{fontenc}
\usepackage{amsmath,amssymb,amsthm}
\usepackage[margin=1in]{geometry}
\usepackage[hidelinks]{hyperref}
\hypersetup{
  pdftitle={The Smallest Eberlein Eigenvalue of J(4r+2,2r,r): The Sharp Threshold r0=2},
  pdfauthor={Xiaojian Zeng}
}

\newtheorem{theorem}{Theorem}
\newtheorem{lemma}[theorem]{Lemma}
\newcommand{\binomzero}[2]{\binom{#1}{#2}}

\title{The Smallest Eberlein Eigenvalue of $J(4r+2,2r,r)$:\\
The Sharp Threshold $r_0=2$}
\author{Xiaojian Zeng\\
PBC School of Finance, Tsinghua University, 100190\\
\texttt{xiaojianzeng@foxmail.com}}
\date{}

\begin{document}
\maketitle

\begin{abstract}
For the generalized Johnson graph $J(4r+2,2r,r)$, we prove that the
Eberlein eigenvalue with index $2$ is the unique smallest adjacency
eigenvalue for every integer $r\geq2$.  Consequently, the least threshold in
the stated domain is $r_0=2$, and the requested exception range is empty.
The proof makes an earlier asymptotic trace argument effective: an exact
multiplicity bound and a strictly decreasing rational trace ratio exclude all
indices at least $4$, while the remaining indices and the case $r=2$ are
handled by direct exact comparisons.
\end{abstract}

\section{Statement and context}

For integers $n,k,s$ with $0\leq s\leq k\leq n$, the generalized Johnson
graph $J(n,k,s)$ has the $k$-subsets of $[n]$ as vertices, with two distinct
vertices adjacent when their intersection has size $s$.  These graphs are
relations of the Johnson association scheme; standard background on the
scheme and its eigenspaces is given by Godsil and Meagher
\cite[Chapter~6]{GodsilMeagher2016}.  Smallest eigenvalues of distance graphs
in the Johnson and other classical schemes have also been studied
systematically by Brouwer, Cioabă, Ihringer, and McGinnis
\cite{BrouwerEtAl2018}.

Fix an integer $r\geq2$.  With the convention that a binomial coefficient is
zero when its lower argument is outside its natural range, write
\begin{equation}\label{eq:eberlein}
 E_r(j)=\sum_{\ell=0}^{r}(-1)^\ell
 \binom{j}{\ell}\binom{2r-j}{r-\ell}
 \binom{2r+2-j}{r-\ell},
 \qquad 0\leq j\leq2r.
\end{equation}
These are the adjacency eigenvalues of $J(4r+2,2r,r)$, and their
multiplicities are
\begin{equation}\label{eq:multiplicity}
 m_j=\binom{4r+2}{j}-\binom{4r+2}{j-1}.
\end{equation}
Luo, Cao, Li, and Li record \eqref{eq:eberlein}--\eqref{eq:multiplicity}
and prove that, in the more general family $J(4r+t,2r,r)$ under
$t^2-t<4r-2$, the index-$2$ eigenvalue is smallest for all sufficiently
large $r$ \cite[Eqns.~(2.1)--(2.2) and Proposition~12]{LuoEtAl2026}.
Their result at $t=2$ does not specify a threshold or classify the remaining
parameters.  We determine both, without making a broader priority claim.

Let
\[
 I(r)=\left\{j\in\{0,\ldots,2r\}:E_r(j)
 =\min_{0\leq i\leq2r}E_r(i)\right\}.
\]

\begin{theorem}\label{thm:main}
For every integer $r\geq2$, one has $I(r)=\{2\}$.  Hence the least integer
$r_0\geq2$ for which $E_r(2)$ is a smallest adjacency eigenvalue for every
$r\geq r_0$ is $r_0=2$, and there are no exceptions with
$2\leq r<r_0$.
\end{theorem}

\section{Exact proof}

Put
\[
 v=\binom{4r+2}{2r},\qquad
 \Delta=\binom{2r}{r}\binom{2r+2}{r},\qquad
 \theta_j=\frac{E_r(j)}{\Delta}.
\]
Thus $v$ is the number of vertices and $\Delta$ is the degree.
Direct substitution in \eqref{eq:eberlein} gives
\begin{equation}\label{eq:low-indices}
 \theta_0=1,\qquad
 \theta_1=\frac{1}{2r+2},\qquad
 \theta_2=-\frac{2r^2+1}{(2r-1)(2r+2)(2r+1)},
\end{equation}
and
\begin{equation}\label{eq:theta3}
 \theta_3-\theta_2=\frac{r-1}{2(r+1)(2r+1)}>0.
\end{equation}
It follows that $\theta_2<\theta_j$ for $j=0,1,3$.  We exclude all
remaining indices through a multiplicity estimate and the spectral trace.

\begin{lemma}\label{lem:multiplicity}
If $r\geq3$, then $m_j\geq m_4$ for every $4\leq j\leq2r$.
\end{lemma}

\begin{proof}
Cancellation in \eqref{eq:multiplicity} gives, for $0\leq j<2r$,
\begin{equation}\label{eq:ratio-m}
 \frac{m_{j+1}}{m_j}=
 \frac{(4r+1-2j)(4r+3-j)}{(j+1)(4r+3-2j)}.
\end{equation}
All factors in \eqref{eq:ratio-m} are positive in this range.  The sign of
the ratio minus one is therefore the sign of
\[
 g_r(j)=j^2-(4r+2)j+4r^2+3r.
\]
Since
\[
 g_r(j+1)-g_r(j)=2j-4r-1<0\qquad (j\leq2r),
\]
the sequence $(m_j)_{4\leq j\leq2r}$ first increases and then decreases,
with either phase possibly empty.  Its minimum is consequently attained at
an endpoint.

It remains to compare the endpoints.  Because $2r\geq6$ and binomial
coefficients increase up to the middle,
\[
 \binom{4r+2}{2r}\geq\binom{4r+2}{5}.
\]
Hence
\[
 m_{2r}=\frac{3}{2r+3}\binom{4r+2}{2r}
 \geq \frac{3}{2r+3}\binom{4r+2}{5}\geq m_4.
\]
After division by $\binom{4r+2}{4}$, the last inequality is
\[
 \frac{3(4r-2)}{5(2r+3)}\geq\frac{4r-5}{4r-1}.
\]
Its cross-multiplied numerator difference is
$8r^2-46r+81$, which is positive because its leading coefficient is
positive and its discriminant is $-476$.  Thus both endpoints, and hence
every term between them, are at least $m_4$.
\end{proof}

\begin{lemma}\label{lem:trace-ratio}
For every integer $r\geq3$,
\[
 \frac{v/\Delta}{m_4\theta_2^2}<1.
\]
\end{lemma}

\begin{proof}
Define
\[
 Q(r)=\frac{v/\Delta}{m_4\theta_2^2}.
\]
Exact factorial cancellation yields
\begin{equation}\label{eq:q-formula}
 Q(r)=
 \frac{6(2r-1)^2\binom{4r}{2r}(r-1)!(r+2)!}
 {(4r-5)(2r^2+1)^2\binom{2r}{r}(2r)!}
\end{equation}
and $Q(3)=330/361<1$.  Consecutive terms satisfy
\begin{equation}\label{eq:q-ratio}
 \frac{Q(r+1)}{Q(r)}=
 \frac{r(r+3)(4r-5)(4r+1)(4r+3)(2r^2+1)^2}
 {2(r+1)(2r-1)^2(2r+1)(4r-1)(2r^2+4r+3)^2}.
\end{equation}
The denominator $D(r)$ on the right of \eqref{eq:q-ratio} is positive for
$r\geq1$, and direct expansion gives
\begin{equation}\label{eq:p-polynomial}
 1-\frac{Q(r+1)}{Q(r)}=\frac{P(r)}{D(r)},
\end{equation}
where
\begin{align*}
 P(r)={}&384r^8+2032r^7+1276r^6+76r^5-236r^4\\
       &\quad+136r^3+439r^2+87r-18.
\end{align*}
For $r\geq1$, the two negative terms are dominated by
\[
 384r^8-236r^4-18\geq148r^4-18>0,
\]
and every omitted term of $P(r)$ is positive.  Therefore $P(r)>0$,
$Q(r)$ is strictly decreasing, and $Q(r)<1$ for all $r\geq3$.
\end{proof}

\begin{proof}[Proof of Theorem~\ref{thm:main}]
The graph is simple and $\Delta$-regular, so
$\operatorname{tr}(A^2)=v\Delta$.  Expanding the trace over the Johnson
scheme eigenspaces and dividing by $\Delta^2$ gives the exact identity
\begin{equation}\label{eq:trace}
 \sum_{j=0}^{2r}m_j\theta_j^2=\frac{v}{\Delta}.
\end{equation}
Suppose $r\geq3$ and that some $j\geq4$ satisfies
$\theta_j\leq\theta_2<0$.  Then
$\theta_j^2\geq\theta_2^2$, while Lemma~\ref{lem:multiplicity} and
\eqref{eq:trace} imply
\[
 \frac{v}{\Delta}\geq m_j\theta_j^2
 \geq m_4\theta_2^2.
\]
This contradicts Lemma~\ref{lem:trace-ratio}.  Hence every $j\geq4$ has
$\theta_j>\theta_2$.  Equations \eqref{eq:low-indices} and
\eqref{eq:theta3} handle $j=0,1,3$, so $I(r)=\{2\}$ for every $r\geq3$.

At the sole remaining parameter $r=2$, exact evaluation of
\eqref{eq:eberlein} gives
\[
 (E_2(0),E_2(1),E_2(2),E_2(3),E_2(4))
 =(90,15,-9,-6,6).
\]
Thus $I(2)=\{2\}$ as well.  Since the allowed domain begins at $2$, the
least threshold is $r_0=2$ and the exception range is empty.
\end{proof}

\section*{Computational and AI assistance}

Exact-arithmetic Python scripts were used to check the low-index formulas,
the multiplicity ratio, the identities \eqref{eq:q-formula}--\eqref{eq:p-polynomial},
the case $r=2$, and finite spectra through $r=60$.  None of these finite
computations is a premise of the quantified proof.  This separation between
diagnostic exact computation and the proof itself is also made explicit in
Section~4 of a recent algebraic graph-invariant calculation \cite{Zeng2026}.
OpenAI Codex assisted with literature triage, proof organization, exact-code
checking, and LaTeX drafting.  The author reviewed the displayed identities,
scope, citations, and final manuscript and takes responsibility for them.

\bibliographystyle{plain}
\bibliography{references}

\end{document}
